\documentclass[11pt]{article}

\usepackage[utf8]{inputenc}
\usepackage[T1]{fontenc}
\usepackage{lmodern}
\usepackage{geometry}
\usepackage{amsmath,amssymb}
\usepackage{siunitx}
\usepackage{booktabs}
\usepackage{enumitem}
\usepackage[round]{natbib}
\usepackage{hyperref}
\usepackage{xcolor}

\geometry{letterpaper, left=1in, right=1in, top=1in, bottom=1in}
\hypersetup{hidelinks=true}

\title{Uniform vertical regridding of radiosonde and dropsonde profiles}
\author{Thomas D. DeWitt}
\date{\today}

\begin{document}
\maketitle

\section{Introduction}

Radiosonde and dropsonde profiles report atmospheric variables at irregular vertical intervals that differ between instruments, processing methods, and field campaigns.  To compare profiles across datasets or to compute structure functions at well-defined separation distances, we first place all profiles on a common vertical grid.  The algorithm described here does so by constructing a uniform grid along the altitude axis and averaging all measurements that fall within each cell. If no measurements fall within a given cell, it is marked as missing.

The simplicity of this approach is deliberate. If interpolation was used, artificial smoothness would be introduced where data are interpolated given that commonly-used interpolation schemes assume a smooth function with $H=1$ \citep[cf.][]{lovejoy2007}.

\subsection{IGRA}

The Integrated Global Radiosonde Archive \citep{durre2006} is a quality-assured compilation of radiosonde observations from more than 1500 globally distributed stations, with records extending from the 1960s to the present.  Observations typically include pressure, temperature, dewpoint depression, geopotential height, wind direction, and wind speed.  IGRA integrates data from disparate sources and applies automated quality assurance that removes measurements failing checks for physical implausiblity, internal inconsistency, climatological outliers, and temporal inconsistency.

We consider data from 2025. To reduce the volume of data, we only consider soundings launched every 10th day (e.g. DOY 1, 11, 21, 31, etc.).

\subsection{JOANNE}

The Joint dropsonde Observations of the Atmosphere in tropical North atlaNtic meso-scale Environments (JOANNE) dataset \citep{george2021} contains 1215 dropsondes deployed during the EUREC\textsuperscript{4}A field campaign in January--February 2020.  Dropsondes were released from the German HALO and NOAA WP-3D aircraft over the tropical North Atlantic east of Barbados, with the goal of characterizing the thermodynamic and dynamic environment of trade-wind cumulus clouds.  Sondes were deployed along circular flight patterns (${\sim}\SI{222}{km}$ diameter) to enable estimation of area-averaged mesoscale divergence, vorticity, and vertical velocity.  The dataset provides both individual sounding profiles and a uniformly gridded product at $\SI{10}{m}$ vertical spacing.  \citet{george2021} assign each profile a quality control flag of ``good,'' ``bad,'' or ``ugly'' based on launch detection, profile completeness ($\geq$80\% for all variables), and plausibility of near-surface values.  A multiplicative correction factor of 1.06 was applied to relative humidity from HALO sondes to compensate for a systematic dry bias caused by insufficient reconditioning of the humidity sensor before flight.

Here, only profiles flagged as ``good'' are used.

\subsection{BEACH}

The Barbados and Eastern Atlantic Combined High-altitude (BEACH) dropsonde dataset \citep{gloeckner2025} contains 1191 dropsondes deployed from the HALO aircraft during the ORCESTRA field campaign in August--September 2024.  Soundings were coordinated by the PERCUSION and MAESTRO subcampaigns and targeted the structure and variability of the Atlantic Inter-tropical Convergence Zone (ITCZ).  Circular flight patterns at approximately $\SI{14}{km}$ altitude (${\sim}\SI{260}{km}$ diameter) were flown over both the eastern and western tropical Atlantic, with flight tracks designed to coincide with EarthCARE satellite overpasses and to sample within and across the ITCZ.  Like JOANNE, the dataset includes a gridded product at $\SI{10}{m}$ spacing and derived mesoscale circulation diagnostics.  \citet{gloeckner2025} assign each profile an aggregate quality flag (GOOD, BAD, or UGLY) based on profile sparsity, profile extent above $\SI{8000}{m}$, near-surface coverage, and physical plausibility of surface values.  Seventeen sondes were falsely configured by the manufacturer as minisondes rather than NRD41 (listed in supplement Table~E2 of \citealt{gloeckner2025}), and one sonde (0bd0e322, flight HALO-20240924a) had a factory-reset sensor module producing anomalous temperature and humidity.

Here, only profiles flagged as ``GOOD'' are used, and the 17 misconfigured and 1 factory-reset sondes listed above are excluded.

\subsection{OTREC}

The Organization of Tropical East Pacific Convection (OTREC) dropsonde dataset \citep{vomel2021} contains 648 dropsondes released from the NSF/NCAR Gulfstream-V during 22 research flights between August and October 2019.  The campaign investigated the distribution and day-to-day variability of deep atmospheric convection in the tropical eastern Pacific and Caribbean, with a focus on why higher rainfall rates occur over lower sea surface temperatures in this region.  Soundings were launched in grid patterns with a typical horizontal spacing of $1^\circ$ longitude and $1.2^\circ$ latitude, from a mean flight altitude of $\SI{13.3}{km}$ to the surface.  OTREC was the first field campaign to rely entirely on the NCAR NRD41 dropsonde, with a nominal vertical resolution between 6 and $\SI{12}{m}$.  All soundings in the distributed dataset have passed ASPEN quality control with manual post-processing corrections; erroneous data points are set to missing.  \citet{vomel2021} applied a per-sonde pressure offset correction (median $\SI{0.35}{hPa}$) and a relative humidity time-lag correction for the slow sensor response above ${\sim}\SI{11.5}{km}$.  One sounding (20190925\_154412) exhibits an unexplained warm bias of 1--$\SI{3.7}{K}$ throughout the profile.

Here, the corrected dataset is used as provided, with the exception of the warm-biased sounding noted above, which is excluded.

\subsection{ACTIVATE}

The Aerosol Cloud meTeorology Interactions oVer the western ATlantic Experiment (ACTIVATE) dropsonde dataset \citep{vomel2023activate} contains 801 dropsondes released from NASA Langley's King Air research aircraft during 165 flights between February 2020 and June 2022.  The campaign studied aerosol--cloud interactions over the western North Atlantic, off the coast of the eastern United States and around Bermuda, sampling all seasons across six deployment periods.  The King Air flew at approximately $\SI{9}{km}$ altitude, and resulting dropsonde profiles provide temperature, pressure, relative humidity, and horizontal and vertical winds between the surface and flight level.  \citet{vomel2023activate} processed all soundings through ASPEN and applied corrected launch times for 17 sondes affected by a synchronization error between the onboard computer and GPS time.  Twenty sondes with unreconditioned humidity sensors may exhibit a dry bias.

Here, profiles marked as degraded are excluded.  The corrected launch times are used where applicable, and only temperature, pressure, and wind from the 20 sondes with unreconditioned humidity sensors is retained.

\subsection{Hurricane dropsonde dataset}

The NOAA hurricane GPS dropsonde dataset \citep{wang2015} consists of 13,681 atmospheric profiles collected during NOAA hurricane reconnaissance and surveillance flights between 1996 and 2012, covering 120 tropical cyclones.  Dropsondes were released from NOAA Gulfstream-IV and WP-3D aircraft in the inner and outer core regions of tropical cyclones, typically at $\SI{150}{}$--$\SI{200}{km}$ horizontal intervals.  The dataset uses the NCAR Airborne Vertical Atmospheric Profiling System (AVAPS), which underwent a major redesign in 2008 (AVAPS II) that expanded simultaneous sounding capacity to eight channels and improved GPS wind recovery rates.  \citet{wang2015} processed all profiles through ASPEN and performed visual inspection; the distributed dataset constitutes a single quality-controlled product.

Here, the archive is used as provided.

\subsection{DYNAMO}

The Dynamics of the Madden--Julian Oscillation (DYNAMO) field campaign \citep{ciesielski2014} investigated the initiation and propagation of the MJO over the tropical Indian Ocean during Fall 2011 through Spring 2012.  In addition to the campaign's extensive radiosonde network, 470 GPS dropsondes were released from the NOAA P-3 aircraft during 13 research flights between 9~November and 13~December 2011.  The P-3 was based at Diego Garcia and flew at an average altitude of $\SI{2}{km}$ (maximum ${\sim}\SI{8}{km}$), with missions sampling the central Indian Ocean ($14.5^\circ$S--$1.2^\circ$N, $69.8^\circ$--$80.9^\circ$E) for convective systems, boundary layer structure, and ITCZ variability.  Of the 470 deployed, 469 passed quality control and are included in the final dropsonde archive.  All profiles were processed through ASPEN; the Version~3 release incorporates a temperature-dependent dry bias correction to relative humidity that was applied to all RD94 dropsondes collected from 2010 onward.

Here, 468 profiles from the dropsonde archive are used as provided.

\subsection{SHOUT}

The Sensing Hazards with Operational Unmanned Technology (SHOUT) project \citep{wick2020} deployed 638 GPS dropsondes from NASA's Global Hawk unmanned aircraft (cruise altitude ${\sim}\SI{18}{km}$) during three field campaigns in 2015 and 2016.  The Global Hawk flew 15 missions sampling 6 tropical cyclones and 3 Pacific winter storms, with missions designed to target regions of high model forecast uncertainty identified through adjoint and ensemble sensitivity analyses.  All profiles were processed through ASPEN.

We use the archive as provided.

\subsection{AR Recon}

Atmospheric River Reconnaissance (AR Recon) is a recurring campaign that has deployed dropsondes from NOAA G-IVSP and Air Force Reserve WC-130J aircraft during winter seasons (November--March) since 2016, targeting atmospheric rivers approaching the U.S.\ West Coast with the primary aim of improving forecasts of landfalling events.  \citet{cobb2022} describe the 2021 season, during which 1142 dropsondes were released during 45 flights.  The data used here span the 2018--2026 winter seasons.  Drop locations are guided by adjoint and ensemble sensitivity analyses to target regions where additional observations would most reduce forecast uncertainty.  All profiles were processed through ASPEN quality control.

Here, the archive is used as provided.

\section{Regridding algorithm}

The output grid is a one-dimensional partition of altitude into cells of uniform thickness $\Delta z$.  Cell edges are placed at
\begin{align}
    z_k = k\,\Delta z, \qquad k = 0, 1, \ldots, N, \label{eq:edges}
\end{align}
where $N = \lfloor z_{\max}/\Delta z\rfloor$ is the number of cells.  The $k$-th cell spans the half-open interval $[z_k,\, z_{k+1})$.  Cell-center altitudes, used as the coordinate for all output variables, are
\begin{align}
    \bar{z}_k = (k + \tfrac{1}{2})\,\Delta z. \label{eq:centers}
\end{align}
The grid spacing is set to $\Delta z = \SI{10}{m}$, chosen to be finer than the effective vertical resolution of most modern radiosondes and dropsondes (of order $\SI{200}{m}$) so that no information is lost during regridding.  The grid extends to $z_{\max} = \SI{20}{km}$ for all dropsonde datasets and $z_{\max} = \SI{40}{km}$ for IGRA radiosondes, which routinely ascend into the stratosphere.

For each grid cell, the regridded value of a variable is equal to the arithmetic mean of all measurements whose altitude falls within that cell.  If a cell contains no measurements, the value is recorded as missing, and this procedure is applied independently to each variable.

All dropsonde datasets report geometric (GPS-derived) altitude.  IGRA reports geopotential height, which is converted to geometric altitude via $z = R_e H / (R_e - H)$, where $H$ is geopotential height and $R_e = \SI{6371}{km}$ is the mean Earth radius.  The difference is approximately 0.5\% at $\SI{30}{km}$ (${\sim}\SI{100}{m}$).

The following variables are bin-averaged onto the uniform grid (Tab. \ref{tab:variables}): zonal and meridional wind ($u$, $v$), pressure ($p$), temperature ($T$), and relative humidity ($\mathrm{RH}$).  When a dataset provides the water vapor mixing ratio $q$ directly rather than relative humidity, $q$ is bin-averaged in place of $\mathrm{RH}$ and no subsequent diagnosis is needed. 


Five additional variables are diagnosed from the regridded fields (Tab. \ref{tab:variables}).  Diagnosis is performed after bin averaging so that the thermodynamic relations are evaluated on the gridded values rather than on individual high-frequency measurements. 

\begin{table}[h]
\centering
\caption{Output variables.  Reported variables are bin-averaged directly; diagnosed variables are computed from the gridded fields after bin averaging.}
\label{tab:variables}
\begin{tabular}{lllll}
\toprule
Variable & Description & Units & Type \\
\midrule
\texttt{u}                  & Zonal wind component             & \si{m\per s}    & Reported \\
\texttt{v}                  & Meridional wind component         & \si{m\per s}    & Reported \\
\texttt{p}                  & Pressure                          & \si{Pa}         & Reported \\
\texttt{T}                  & Temperature                       & \si{K}          & Reported \\
\texttt{RH}                 & Relative humidity                 & \%              & Reported \\
\texttt{observation\_time}  & Measurement time at each level    & datetime        & Reported \\
\texttt{q}                  & Water vapor mixing ratio          & \si{kg\per kg}  & Diagnosed (Eq.~\ref{eq:q}) \\
\texttt{theta}              & Potential temperature              & \si{K}          & Diagnosed (Eq.~\ref{eq:theta}) \\
\texttt{theta\_e}           & Equivalent potential temperature   & \si{K}          & Diagnosed (Eq.~\ref{eq:theta_e}) \\
\texttt{MSE}                & Moist static energy               & \si{J\per kg}   & Diagnosed (Eq.~\ref{eq:MSE}) \\
\texttt{DSE}                & Dry static energy                 & \si{J\per kg}   & Diagnosed (Eq.~\ref{eq:DSE}) \\
\bottomrule
\end{tabular}
\end{table}

When relative humidity is the reported moisture variable, the mixing ratio is computed from the regridded $\mathrm{RH}$, $T$, and $p$ fields.  The saturation vapor pressure over liquid water $e_s$ is obtained from the August--Roche--Magnus formula,
\begin{align}
    e_s(T) = 610.94 \exp\!\left(\frac{17.625\, T_c}{T_c + 243.04}\right), \label{eq:es}
\end{align}
where $T_c = T - \SI{273.15}{K}$ is the temperature in degrees Celsius and $e_s$ has units of Pascals.  The vapor pressure is $e = (\mathrm{RH}/100)\, e_s$, and the mixing ratio follows as
\begin{align}
    q = \frac{0.622\, e}{p - e}. \label{eq:q}
\end{align}


Potential temperature is computed from Poisson's equation,
\begin{align}
    \theta = T \left(\frac{p_0}{p}\right)^{\kappa}, \label{eq:theta}
\end{align}
where $p_0 = \SI{100000}{Pa}$ is the reference pressure and $\kappa = R_d / c_p \approx 0.2854$ is the Poisson constant, with $R_d = \SI{287.04}{J\,kg^{-1}\,K^{-1}}$ the specific gas constant for dry air and $c_p = \SI{1005.7}{J\,kg^{-1}\,K^{-1}}$ the isobaric specific heat.


Moist static energy combines sensible heat, gravitational potential energy, and latent heat into a single conserved quantity for moist adiabatic processes:
\begin{align}
    h = c_p T + g\bar{z}_k + L_v q, \label{eq:MSE}
\end{align}
where $g = \SI{9.807}{m\,s^{-2}}$ is gravitational acceleration, $L_v = \SI{2.501e6}{J\,kg^{-1}}$ is the latent heat of vaporization at $\SI{0}{\celsius}$, and $\bar{z}_k$ is the cell-center altitude from Eq.~\ref{eq:centers}.


Dry static energy is conserved for dry adiabatic processes:
\begin{align}
    s = c_p T + g\bar{z}_k. \label{eq:DSE}
\end{align}


Equivalent potential temperature is conserved for saturated adiabatic processes.  It is approximated following \citet{bolton1980}:
\begin{align}
    \theta_e = \theta \exp\!\left(\frac{L_v\, q}{c_p\, T_L}\right), \label{eq:theta_e}
\end{align}
where $T_L$ is the temperature at the lifting condensation level, estimated as
\begin{align}
    T_L = \frac{1}{\displaystyle\frac{1}{T - 55} - \frac{\ln(\mathrm{RH}/100)}{2840}} + 55. \label{eq:TL}
\end{align}

\subsection{Output format} \label{sec:output}

The regridded data are stored as one NetCDF-4 file per source dataset, or one file per year for IGRA.  Each file has three dimensions: \texttt{launch\_location}, \texttt{sounding}, and \texttt{altitude}.
For radiosonde datasets, every launch from a given station shares the \texttt{launch\_location} coordinate. The size of the \texttt{sounding} dimension is 1 for dropsonde datasets given that each dropsonde is launched from a different spatial location.
The \texttt{altitude} coordinate contains the cell-center altitudes $\bar{z}_k$ from Eq.~\ref{eq:centers}, in meters above mean sea level.

Each launch location carries the following metadata:

\begin{center}
\begin{tabular}{lll}
\toprule
Variable & Type & Description \\
\midrule
\texttt{launch\_time}        & datetime & Date and time of sounding launch (\texttt{launch\_location}, \texttt{sounding}) \\
\texttt{observation\_time}  & datetime & Observation time at each altitude level (\texttt{launch\_location}, \texttt{sounding}, \texttt{altitude}) \\
\texttt{launch\_lat}        & float    & Latitude at profile start [\si{\degree}N] \\
\texttt{launch\_lon}        & float    & Longitude at profile start [\si{\degree}E] \\
\texttt{sonde\_id}          & string   & Provider's sonde or profile identifier (dropsonde datasets) \\
\texttt{station\_id}        & string   & Station identifier (radiosonde datasets) \\
\bottomrule
\end{tabular}
\end{center}

\noindent The \texttt{observation\_time} variable records the actual measurement time at each altitude level, bin-averaged within each grid cell.  Each profile has either \texttt{sonde\_id} or \texttt{station\_id}, not both.  For dropsonde datasets, \texttt{launch\_lat} and \texttt{launch\_lon} refer to the aircraft position at the time of release, or the first GPS measurement in the profile when explicit release coordinates are not available.  For radiosondes, they refer to the launch site.

\subsection{Data variables}

All regridded and diagnosed variables are stored with shape (\texttt{launch\_location}, \texttt{sounding}, \texttt{altitude}).  Missing values (grid cells with no measurements) are stored as NaN.



\section{Validation} \label{sec:validation}

Two of the dropsonde datasets---JOANNE \citep{george2021} and BEACH \citep{gloeckner2025}---provide both individual sounding profiles and a uniformly gridded product at $\SI{10}{m}$ vertical spacing produced by the data providers. 
Agreement between the two gridded products is expected to be close but not exact given that the provider gridded products use interpolation rather than bin averaging, and in the case of BEACH use slightly different filtering criteria.

For each variable, the comparison is performed per sounding at all grid cells where both products report non-missing values.  Across all cells, the median RMSD and mean bias between our regridded product and the official products are negligible (Tab. \ref{tab:validation}), with the possible exception of the BEACH pressure variable. A small discrepancy is to be expected here given that the gridded pressure variable from \citet{gloeckner2025} obtained pressure through integrating the thermodynamic profile, not interpolating measured pressure values.

\begin{table}[h]
\centering
\caption{Validation statistics comparing bin-averaged profiles to provider-gridded products.  RMSD and bias are computed per sounding over matched non-missing grid cells; the table reports the median RMSD and mean bias across all matched sondes.}
\label{tab:validation}
\begin{tabular}{llrrr}
\toprule
Dataset & Variable & Median RMSD & Mean bias & Units \\
\midrule
JOANNE (1068 sondes) & $T$             & 0.012  & $<$0.001      & \si{K} \\
                      & $u$             & 0.025  & $-$0.000      & \si{m\per s} \\
                      & $v$             & 0.024  & $<$0.001      & \si{m\per s} \\
                      & $\mathrm{RH}$   & 0.12   & $<$0.001      & \% \\
                      & $p$             & 11.3   & $-$0.003      & \si{Pa} \\
\midrule
BEACH (976 sondes)    & $T$             & 0.24   & $-$0.18       & \si{K} \\
                      & $u$             & 0.27   & $-$0.023      & \si{m\per s} \\
                      & $v$             & 0.26   & $<$0.001      & \si{m\per s} \\
                      & $\mathrm{RH}$   & 1.65   & $-$0.091      & \% \\
                      & $p$             & 137    & $-$120        & \si{Pa} \\
\bottomrule
\end{tabular}
\end{table}

\bibliographystyle{apalike}
\bibliography{references}

\end{document}
